Bodový graf závislosti průměrného počtu odpracovaných hodin za rok na pokrytí \ac{KS} ukazuje zápornou korelaci: pracovníci v~zemích s~vyšším \ac{KS} pokrytím pracují v~průměru méně hodin ročně při srovnatelném nebo vyšším reálném příjmu --- efekt, který je konzistentní s~teorií, podle níž silné \ac{KS} prosazují i~pracovní dobu jako vyjednávací proměnnou (zkrácené pracovní úvazky, limit přesčasů).
V~ČR, kde \ac{KS} zřídka obsahují robustní ustanovení o~pracovní době nad zákonné minimum, pracovníci ve srovnatelných odvětvích odpracují více hodin než v~AT nebo DK, přičemž hodinová mzdová sazba je přesto nižší --- kombinace více hodin za méně peněz je symptomem absence koordinovaného vyjednávání.
Prodloužená pracovní doba v~ČR (nadprůměrný podíl přesčasů, nejnižší podíl částečných úvazků v~\ac{EU}) nesignalizuje vyšší pracovní nasazení, ale neschopnost kolektivní smlouvy prosadit kompenzaci produktivity formou kratší doby nebo vyšší sazby za přesčas.
Reforma \ac{KV} by proto měla výslovně zahrnout pracovní dobu jako standardní vyjednávací agendu: \ac{KS}, která řeší jen základní mzdu bez doby, nevyjednává celkovou hodnotu pracovního vztahu.

Bodový graf závislosti průměrného počtu odpracovaných hodin za rok na pokrytí \ac{KS} ukazuje zápornou korelaci: pracovníci v~zemích s~vyšším \ac{KS} pokrytím pracují v~průměru méně hodin ročně při srovnatelném nebo vyšším reálném příjmu --- efekt, který je konzistentní s~teorií, podle níž silné \ac{KS} prosazují i~pracovní dobu jako vyjednávací proměnnou (zkrácené pracovní úvazky, limit přesčasů).
V~ČR, kde \ac{KS} zřídka obsahují robustní ustanovení o~pracovní době nad zákonné minimum, pracovníci ve srovnatelných odvětvích odpracují více hodin než v~AT nebo DK, přičemž hodinová mzdová sazba je přesto nižší --- kombinace více hodin za méně peněz je symptomem absence koordinovaného vyjednávání.
Prodloužená pracovní doba v~ČR (nadprůměrný podíl přesčasů, nejnižší podíl částečných úvazků v~\ac{EU}) nesignalizuje vyšší pracovní nasazení, ale neschopnost kolektivní smlouvy prosadit kompenzaci produktivity formou kratší doby nebo vyšší sazby za přesčas.
Reforma \ac{KV} by proto měla výslovně zahrnout pracovní dobu jako standardní vyjednávací agendu: \ac{KS}, která řeší jen základní mzdu bez doby, nevyjednává celkovou hodnotu pracovního vztahu.

Bodový graf závislosti průměrného počtu odpracovaných hodin za rok na pokrytí \ac{KS} ukazuje zápornou korelaci: pracovníci v~zemích s~vyšším \ac{KS} pokrytím pracují v~průměru méně hodin ročně při srovnatelném nebo vyšším reálném příjmu --- efekt, který je konzistentní s~teorií, podle níž silné \ac{KS} prosazují i~pracovní dobu jako vyjednávací proměnnou (zkrácené pracovní úvazky, limit přesčasů).
V~ČR, kde \ac{KS} zřídka obsahují robustní ustanovení o~pracovní době nad zákonné minimum, pracovníci ve srovnatelných odvětvích odpracují více hodin než v~AT nebo DK, přičemž hodinová mzdová sazba je přesto nižší --- kombinace více hodin za méně peněz je symptomem absence koordinovaného vyjednávání.
Prodloužená pracovní doba v~ČR (nadprůměrný podíl přesčasů, nejnižší podíl částečných úvazků v~\ac{EU}) nesignalizuje vyšší pracovní nasazení, ale neschopnost kolektivní smlouvy prosadit kompenzaci produktivity formou kratší doby nebo vyšší sazby za přesčas.
Reforma \ac{KV} by proto měla výslovně zahrnout pracovní dobu jako standardní vyjednávací agendu: \ac{KS}, která řeší jen základní mzdu bez doby, nevyjednává celkovou hodnotu pracovního vztahu.

Bodový graf závislosti průměrného počtu odpracovaných hodin za rok na pokrytí \ac{KS} ukazuje zápornou korelaci: pracovníci v~zemích s~vyšším \ac{KS} pokrytím pracují v~průměru méně hodin ročně při srovnatelném nebo vyšším reálném příjmu --- efekt, který je konzistentní s~teorií, podle níž silné \ac{KS} prosazují i~pracovní dobu jako vyjednávací proměnnou (zkrácené pracovní úvazky, limit přesčasů).
V~ČR, kde \ac{KS} zřídka obsahují robustní ustanovení o~pracovní době nad zákonné minimum, pracovníci ve srovnatelných odvětvích odpracují více hodin než v~AT nebo DK, přičemž hodinová mzdová sazba je přesto nižší --- kombinace více hodin za méně peněz je symptomem absence koordinovaného vyjednávání.
Prodloužená pracovní doba v~ČR (nadprůměrný podíl přesčasů, nejnižší podíl částečných úvazků v~\ac{EU}) nesignalizuje vyšší pracovní nasazení, ale neschopnost kolektivní smlouvy prosadit kompenzaci produktivity formou kratší doby nebo vyšší sazby za přesčas.
Reforma \ac{KV} by proto měla výslovně zahrnout pracovní dobu jako standardní vyjednávací agendu: \ac{KS}, která řeší jen základní mzdu bez doby, nevyjednává celkovou hodnotu pracovního vztahu.

\input{texparts/python/coverage_hours_scatter}



